% Final §4 -- Evaluation (plan §7.1: 2.5 pages). Every measured value is a macro of
% generated/numbers.tex (source file and column in its comments and in generated/README.md);
% charts are report/figures/generated/*.pdf from report/tools/make_figures.py.
\section{Evaluation}\label{sec:evaluation}

\subsection{Method}
All system experiments ran against the \code{dev} stage through the public API, exactly as the
browser does: client \code{commandId}, POST, 202, then polling with the back-off of
\S\ref{sec:timers}. Each script asserts its checks, calls \code{verify} after its scenario steps
and writes \code{results/<date>-<experiment>/}; each experiment ran 3 times in one invocation.
Expiry experiments use $\text{holdSeconds} = 60$ so that a sample does not take 15 minutes. The
acceptance smoke test passed before the series (\SmokeChecksPassed/\SmokeChecksTotal{} checks, run
\code{\SmokeRunId}). The harness recorded the local commit with uncommitted changes in the working
tree (\code{commit.dirty} in \code{meta.json}):
\pending{confirm the deployed commit and that only results/ and docs/ were uncommitted, or re-run
from a clean checkout}. Run ids and commits of every number are listed in
\code{report/generated/numbers.tex}.

\subsection{System results against the Table~3 targets}\label{sec:results}
Table~\ref{tab:results} lists one row per Table~3 target of the Preliminary report.

\begin{table}[t]
\centering\footnotesize
\caption{Table~3 targets vs measured values (dev stage, 3 runs per experiment unless stated). Source
directories: \code{results/2026-10-17-<experiment>/}, \code{results/replay/}, \code{results/cost/}.}
\label{tab:results}
\begin{tabularx}{\textwidth}{@{}P{1.95cm}P{4.3cm}LP{1.25cm}@{}}
\toprule
Experiment & Target & Measured & Verdict\\
\midrule
Core correctness & 0 over-capacity, 0 double holds, every step as expected, \code{verify} ok
 & \CorrectnessScenarios{} scenarios $\times$ \CorrectnessRuns{} runs: \CorrectnessChecksPassed/\CorrectnessChecksTotal{} checks; over-capacity \CorrectnessOverCapacity, double holds \CorrectnessDoubleHolds; \CorrectnessVerifications{} verify calls, \CorrectnessVerifyFailures{} failed
 & \CorrectnessVerdict\\
Preference, order & Never offer an unacceptable session; preference order respected (engine property test and 20 cases on dev)
 & Engine: \IfPending{\CorePropertySequences}{\pending{core property-test run not yet recorded}}{\CorePropertySequences{} random sequences, \CoreTestsPassed/\CoreTestsTotal{} tests passed, branch coverage \CoreBranchCoveragePct\,\%}. Dev: \PrefOrderCasesPassed/\PrefOrderCases{} cases (\PrefOrderDistinctCases{} $\times$ \PrefOrderRuns{} runs), \PrefOrderNonAcceptable{} unacceptable offers
 & \PrefOrderVerdict{} (dev)\\
Concurrency & Invariants hold after each concurrent phase; version-conflict retries recorded (expected 0)
 & \ConcurrencyAccounts{} accounts race for \ConcurrencySeats{} seats in \ConcurrencySessions{} sessions: \code{verify} found no violation after any phase; retries \ConcurrencyMetaRetries{} (inconclusive, see text); \ConcurrencyTransportErrors{} of \ConcurrencyCommands{} requests got HTTP 503 (see text); \ConcurrencyChecksPassed/\ConcurrencyChecksTotal{} checks
 & \ConcurrencyVerdict\\
Repeat execution & Same \code{commandId} sent repeatedly and crash after commit: applied once, one record
 & \RepeatProbesPassed/\RepeatProbes{} probes with \RepeatCopies{} concurrent copies each; \RepeatAfterCommitPassed/\RepeatAfterCommitCases{} crash-after-commit redeliveries acknowledged without a second effect
 & \RepeatVerdict\\
Expiry race & ACCEPT at \code{expiresAt}$-$200\,ms applied, at $+$200\,ms rejected OFFER\_EXPIRED whatever the timer did; a stale tick has no effect
 & Early applied \ExpiryRaceEarlyApplied/\ExpiryRaceEarlyAttempts{} (see text); late rejected \ExpiryRaceLateRejected/\ExpiryRaceLateTrials; stale ticks without effect \ExpiryRaceLateTicksNoEffect/\ExpiryRaceLateTicks
 & \ExpiryRaceVerdict\\
Failure recovery & Crashes before commit, while sending timers and after commit recover, settle once and are traceable; notifier backlog delivered $<$\,300\,s after an outage
 & Before commit \FailureBeforeCommitPassed/\FailureBeforeCommitCases{} and timer send \FailureTimerSendPassed/\FailureTimerSendCases{} settled on the redelivery (median \FailureBeforeCommitSettleMedianS{} and \FailureTimerSendSettleMedianS\,s); after commit \FailureAfterCommitPassed/\FailureAfterCommitCases{} settled at once (median \FailureAfterCommitSettleMedianS\,s), redelivery acknowledged; backlog of \FailureNotifierOutageNotifications{} notifications drained $\leq$\,\FailureNotifierOutageMaxS\,s after restore
 & \FailureVerdict\\
Access control & Unauthorised requests refused (401/403/404, OFFER\_NOT\_YOURS); nothing changed
 & \IfPending{\AccessVerdict}{\pending{access-control run (dev) not yet complete}}{\AccessProbesPassed/\AccessProbes{} probes in \AccessProbeGroups{} groups as expected}
 & \IfPending{\AccessVerdict}{--}{\AccessVerdict}\\
Performance & Read and 202 p95 $<$\,1\,s; settle p95 $<$\,5\,s (server and client); errors $<$\,1\,\%
 & \IfPending{\PerfDevHotspotVUFiftyServerSettlePNinetyFiveMs}{\pending{k6 matrix (dev) not yet complete}}{50 VUs, cross-workshop: read \PerfDevCrossVUFiftyReadPNinetyFiveMs, 202 \PerfDevCrossVUFiftyAcceptPNinetyFiveMs, settle \PerfDevCrossVUFiftyServerSettlePNinetyFiveMs{} (server) / \PerfDevCrossVUFiftySettlePNinetyFiveMs{} (client) ms; errors \PerfDevCrossVUFiftyErrorPct\,\%. Hot spot: settle \PerfDevHotspotVUFiftyServerSettlePNinetyFiveMs{} ms (server)}
 & Fig.~\ref{fig:perf}\\
Backfill latency & Expiry-driven p95 $\leq$\,120\,s; every expiry relayed; \code{SweepLateExpiries} = 0
 & \IfPending{\BackfillExpiredPNinetyFiveS}{\IfPending{\BackfillVerdict}{\pending{backfill run (dev) not yet complete}}{Run aborted: HTTP 503 when creating its workshop (see text); \pending{backfill re-run}}}{p50 / p95 / max \BackfillExpiredPFiftyS{} / \BackfillExpiredPNinetyFiveS{} / \BackfillExpiredMaxS\,s (n = \BackfillExpiredN), \BackfillExpiredUnrelayed{} not relayed; cancellation-driven max \BackfillCancelMaxMs\,ms; \code{SweepLateExpiries} \BackfillSweepLateExpiries}
 & \IfPending{\BackfillVerdict}{--}{\BackfillVerdict}\\
Cost & Unit cost and trend by scale in three layers; on-prem comparison
 & Model: \usd\CostUnitOneKCloudListUsd{} $\to$ \usd\CostUnitHundredKCloudListUsd{} per 1{,}000 registrations at list price (1k $\to$ 100k a month); crossovers in \S\ref{sec:business}. Measured usage: \CostMeasuredListUsd
 & model\\
Mechanism comparison & Duplicate holds, distinct participants served and refill time, UNIFIED vs PER\_SESSION, 95\,\% CIs
 & \ReplaySeeds{} simulated workshops: duplicate-hold events \ReplayDupHoldEventsUnifiedMean{} vs \ReplayDupHoldEventsPerSessionMean; distinct served, per-session $-$ unified: \ReplayDistinctDeltaMean{} [\ReplayDistinctDeltaLo, \ReplayDistinctDeltaHi]
 & simulation\\
User study & Completion $\geq$\,90\,\%; paired task time $-$25\,\%; System Usability Scale (SUS) $\geq$\,75; 12 participants
 & \IfPending{\StudyParticipants}{\pending{user study not yet run}}{n = \StudyParticipants; completion \StudyCompletionSRPct\,\% (control \StudyCompletionCtrlPct\,\%); time $-$\StudyTimeImprovementPct\,\% (Wilcoxon $p$ = \StudyTimeWilcoxonP); SUS \StudySusSRMean{} [\StudySusSRLo, \StudySusSRHi]}
 & \S\ref{sec:study}\\
Real pilot & Two real sessions run with SeatRelay
 & \IfPending{\PilotWorkshops}{\pending{pilot not yet run}}{\PilotWorkshops{} workshops, \PilotRegistrations{} registrations, \PilotConfirmed{} confirmed, \PilotBackfills{} backfills, \PilotAttended{} attended}
 & \S\ref{sec:pilot}\\
\bottomrule
\end{tabularx}
\end{table}

\paragraph{Failed checks.}
Three experiments failed a check, and we report them as failures. In the \emph{concurrency} run,
\ConcurrencyTransportErrors{} of \ConcurrencyCommands{} requests received HTTP 503 responses that
carry no application error code, i.e.\ they did not come from the API's own 503 path (failed
enqueue), which returns a code. Because some of these commands were nevertheless applied, the
harness lost track of them: its expected sequence numbers and confirmation counts no longer matched,
although \code{verify} found no over-capacity, double hold or other invariant violation after any
phase. The version-conflict metric read \ConcurrencyMetaRetries, but the CloudWatch query returned no
samples for any SeatRelay metric in that window, so that reading is inconclusive.
The \emph{backfill} run aborted at its first request, creating a workshop, with the same kind of
HTTP 503 response, so it produced no latency samples.
\pending{root cause of the HTTP 503 responses, then clean concurrency and backfill re-runs}. In the \emph{expiry
race}, \ExpiryRaceEarlyApplied{} of \ExpiryRaceEarlyAttempts{} early ACCEPTs were applied; the other
reached the API after \code{expiresAt} because of client send-timing jitter, and the server rejected
it as expired, which agrees with the server clock and with C5. The check still counts it as a
failure, and the experiment should be repeated with a larger timing margin. In all
\ExpiryRaceOrderTimerFirst{} late trials the timer's tick had expired the offer before the ACCEPT
arrived, and no stale tick changed a newer offer.

\paragraph{Latency.}
In the correctness run the server settled commands with p95 \CorrectnessSettleServerPNinetyFiveMs\,ms
(\CorrectnessSettleClientN{} commands); the client observed p95 \CorrectnessSettleClientPNinetyFiveMs\,ms,
because it adds the first 300\,ms poll interval, the request round trips and up to one further
interval. In the concurrency run, where all commands of one workshop are serialised in one FIFO
group by design, the server p95 rose to \ConcurrencySettleServerPNinetyFiveMs\,ms. After a worker
crash a command settles when the 120\,s visibility timeout returns the message (median
\FailureBeforeCommitSettleMedianS\,s), which bounds the recovery of a stalled workshop.
\ifloadtest Fig.~\ref{fig:perf} shows the k6 load test~\cite{k6} with 10--50 virtual users (VUs),
for load spread over 10 workshops and for one hot workshop.
\else Fig.~\ref{fig:perf} shows the settle-latency distributions of the functional runs; the k6 load
test~\cite{k6} (10--50 virtual users over 10 workshops and on one hot workshop) is not complete yet.\fi{}
The backfill experiment lets holds expire behind waiting participants and measures the relay through
the timer path (Fig.~\ref{fig:backfill}).

\begin{figure}[t]
  \centering
  \ifloadtest
  \resultfig{loadtest-p95}{5.6cm}{k6 load-test chart}
  \caption{Load test on dev: p95 latency of reads, 202 responses and settlement (server and client)
  by number of virtual users, for load spread over 10 workshops and for one hot workshop (median of
  3 rounds per level). Source: \code{results/loadtest/<date>/}.}
  \else
  \resultfig{settle-latency}{5.6cm}{settle-latency chart}
  \caption{Command settle latency in the functional experiments on dev (empirical cumulative distributions, log scale):
  (a)~server side, \code{settledAt} $-$ \code{acceptedAt}; (b)~client-observed, from the POST to the
  poll that saw the terminal state. Dotted: 95\,\%. Sources:
  \code{results/2026-10-17-\{smoke,correctness,preference-order,concurrency\}/}.}
  \fi
  \label{fig:perf}
\end{figure}

\begin{figure}[t]
  \centering\resultfig{backfill-latency}{5.2cm}{backfill-latency chart (results/<date>-backfill)}
  \caption{Expiry-driven backfill latency (\code{reallocatedAt} $-$ \code{expiresAt}) on dev with
  $\text{holdSeconds} = 60$, histogram and empirical cumulative distribution against the 120\,s target. Source:
  \code{results/<date>-backfill/raw.csv}.}
  \label{fig:backfill}
\end{figure}

\subsection{Mechanism comparison: one queue vs per-session forms}\label{sec:mechanism}
The replay tool runs the production engine with a simulated clock and ideal timers on
\ReplaySeeds{} random workshops (2--5 sessions, \ReplayParamOversubscriptionMin--\ReplayParamOversubscriptionMax{}
times oversubscribed, holds of 60--900\,s, random accept, decline, ignore, cancel and hoarding
behaviour). Each seed's identical arrival stream is replayed twice: in UNIFIED mode (SeatRelay) and in
PER\_SESSION mode, where a participant registers separately for every acceptable session, each
session has its own queue, and duplicates are released by hand only with some probability. With one
queue there were no duplicate holds, as C2 requires; with per-session queues there were
\ReplayDupHoldEventsPerSessionMean{} duplicate-hold events per workshop
[\ReplayDupHoldEventsPerSessionLo, \ReplayDupHoldEventsPerSessionHi] (Fig.~\ref{fig:replay}). One
queue served \ReplayDistinctUnifiedMean{} distinct participants per workshop against
\ReplayDistinctPerSessionMean{} (paired difference \ReplayDistinctDeltaMean{}
[\ReplayDistinctDeltaLo, \ReplayDistinctDeltaHi]; fewer under PER\_SESSION in
\ReplayDistinctFewerSeedsPct\,\% of seeds, equal in \ReplayDistinctEqualSeedsPct\,\%), and a released
seat reached a new participant after \ReplayRefillMeanUnifiedMean{} instead of
\ReplayRefillMeanPerSessionMean{} minutes on average (\ReplayRefillMeanN{} seeds with refills in both
modes). Seats left empty at cut-off did not differ (difference \ReplaySeatsWastedDeltaMean{}
[\ReplaySeatsWastedDeltaLo, \ReplaySeatsWastedDeltaHi]). The comparison is a simulation with an
assumed behaviour model, not a field measurement; the plan's additional replay on the deployed system
was not carried out.

\begin{figure}[t]
  \centering\resultfig{replay-comparison}{5.2cm}{replay chart}
  \caption{Mechanism comparison over \ReplaySeeds{} simulated workshops (one box per mode: p25--p75,
  median line, whiskers p5--p95, diamond = mean). (a)~Duplicate-hold events per workshop;
  (b)~distinct participants served as a share of capacity; (c)~mean seat refill time per workshop
  (log scale). Source: \code{results/replay/raw.csv}.}
  \label{fig:replay}
\end{figure}

\subsection{User study}\label{sec:study}
The protocol has twelve people outside the team (at least four with organising experience) use
SeatRelay and a control that mirrors current practice, one Google Form per session plus a shared Sheet with
duplicate-detection formulas, in counterbalanced order with equivalent data sets. Participants
perform five tasks (choose acceptable sessions, join the waitlist, interpret their status, confirm,
cancel) and organisers four (create a multi-session workshop, handle a cancellation, inspect the
backfill, read the distinct count and export). We measure completion, paired task time (Wilcoxon
signed-rank test~\cite{wilcoxon1945}) and SUS~\cite{brooke1996sus}.
Participants: \StudyParticipants. Completion: \StudyCompletionSRPct\,\% with SeatRelay. Median
paired time improvement: \StudyTimeImprovementPct\,\%. SUS: \StudySusSRMean.

\begin{figure}[t]
  \centering
  \IfFileExists{figures/generated/study-time.pdf}{%
    \includegraphics[width=0.49\textwidth]{figures/generated/study-time}\hfill
    \includegraphics[width=0.49\textwidth]{figures/generated/study-sus}}{%
    \fbox{\parbox[c][\dimexpr6.0cm-2\fboxsep-2\fboxrule\relax][c]{\dimexpr\textwidth-2\fboxsep-2\fboxrule\relax}{%
      \centering\small\pending{user-study charts: paired task time and SUS (results/user-study/)}\par}}}
  \caption{User study: total task time per participant in both conditions (left) and SUS scores
  (right), SeatRelay vs one form per session plus a Sheet.}
  \label{fig:study}
\end{figure}

\subsection{Real pilot}\label{sec:pilot}
The \code{pilot} stage (fault injection disabled, e-mail enabled) serves a real workshop in two
sessions with registration open for at least five days; the organiser's CSV exports and the
dashboard's backfill log are the evidence. \pending{pilot workshop, dates and observations: natural
cancellations, backfills, notifications, attendance, organiser feedback}.
